\subsection{代数的张量积上的半单模}

在本小节中，我们考虑 K-代数 \(A_{1}\) 与 \(A_{2}\)；我们把代数 \(A_{1} \otimes A_{2}\) 记作 A。

命题 1. —— 设 \(M_{1}\) 是一个 \(A_{1}\)-模，\(M_{2}\) 是一个 \(A_{2}\)-模，二者都不约化为 0。若模 \(M = M_{1} \otimes M_{2}\)（环 \(A = A_{1} \otimes A_{2}\) 上的模）是单模（相应地，同构模、半单模），则 \(A_{1}\)-模 \(M_{1}\) 与 \(A_{2}\)-模 \(M_{2}\) 是单模（相应地，同构模、半单模）。

假设 M 是半单 A-模。设 \(N_{1}\) 是一个 \(A_{1}\)-子模（\(M_{1}\) 的）。用 N 表示 \(N_{1} \otimes M_{2}\) 在 A-模 M 中的典范像。由 VIII, 第 56 页，推论 2，存在 M 中一个以 N 为像的 A-线性投影子 p。由假设，\(M_{2} \neq 0\)；因此可以选取 \(M_{2}\) 的一个元素 m 和一个线性形式 \(\varphi\)（在 K-向量空间 \(M_{2}\) 上，满足 \(\varphi(m) = 1\)）（II, §7, No.~5, 第 302 页，推论 2）。设 u 是从 \(M_{1}\) 到 M 的映射，由 \(u(m_{1}) = m_{1} \otimes m\) 定义；v 是从 M 到 \(M_{1}\) 的 K-线性映射，由 \(v(m_{1} \otimes m_{2}) = \varphi(m_{2})m_{1}\) 刻画。令 \(q = v \circ p \circ u\)。映射 \(q : M_{1} \to M_{1}\) 是 \(A_{1}\)-线性的，它的像含于 \(N_{1}\)，并且我们有 \(q(n) = n\)（对每个 \(n \in N_{1}\)）。于是，q 是 \(M_{1}\) 中以 \(N_{1}\) 为像的投影子。我们证明了 \(M_{1}\) 是半单 \(A_{1}\)-模（VIII, 第 56 页，推论 2）。

假设 M 是单模，且 \(M_{1}\) 是两个 \(A_{1}\)-子模 \(M_{1}^{\prime}\) 与 \(M_{1}^{\prime\prime}\) 的直和。令 \(M^{\prime}=M_{1}^{\prime}\otimes M_{2}\) 与 \(M^{\prime\prime}=M_{1}^{\prime\prime}\otimes M_{2}\)；则 M 是 A-子模 \(M'\) 与 \(M''\) 的直和。由于 M 是单模，\(M'\) 或 \(M''\) 必约化为 0；由假设 \(M_{2} \neq 0\)，故 \(M_{1}' = 0\) 或 \(M_{1}'' = 0\)（II, §3, No.~7, 第 256 页，推论 2）。这就证明了 \(M_{1}\) 是单 \(A_{1}\)-模。

现在假设 M 是同构 A-模。设 S 与 T 是两个单 \(A_{1}\)-子模（属于 \(M_{1}\)）。于是 A-模 \(S \otimes M_{2}\) 与 \(T \otimes M_{2}\) 可与 M 的非零子模视为同一。因而它们是与 M 同型的同构模。由注 VIII, 第 61 页，存在一个非零 A-线性映射 \(f : S \otimes M_{2} \to T \otimes M_{2}\)。映射 f 特别还是 \(A_{1}\)-线性的。由于 \(A_{1}\)-模 \(S \otimes M_{2}\) 与 \(T \otimes M_{2}\) 非零，且分别是 S 型与 T 型的同构模，故 S 与 T 同构。这就证明了 \(M_{1}\) 是同构 \(A_{1}\)-模。

命题 2. —— 设 S 是环 \(A = A_{1} \otimes A_{2}\) 上的一个单模，在 K 上有限维。对 \(i \in \{1, 2\}\)，存在一个单 \(A_{i}\)-模 \(S_{i}\)，使得 \(A_{i}\)-模 S 是 \(S_{i}\) 型的同构模。A-模 S 同构于 A-模 \(S_{1} \otimes S_{2}\) 的一个商。

由于 S 是一个在 K 上有限维的非零 \(A_{1}\)-模，它在 \(A_{1}\) 上有有限长度，且存在一个单左 \(A_{1}\)-模 \(S_{1}\) 以及一个非零 \(A_{1}\)-线性映射（从 \(S_{1}\) 到 S）。我们赋予 \(\mathrm{M}_{2} = \mathrm{Hom}_{\mathrm{A}_{1}}(\mathrm{S}_{1}, \mathrm{S})\) 一个左 \(A_{2}\)-模结构，它由外合成律 \((a_{2}, u) \mapsto (a_{2})_{\mathrm{S}} \circ u\) 定义。由构造，\(M_{2} \neq 0\)，且 \(M_{2}\) 在 K 上有限维。因此可以找到一个单左 \(A_{2}\)-模 \(S_{2}\) 与一个非零 \(A_{2}\)-线性映射 \(\varphi : S_{2} \to M_{2}\)。我们定义一个非零 A-线性映射 \(\psi\)（从 \(S_{1} \otimes S_{2}\) 到 S），由

\[
\psi (s _ {1} \otimes s _ {2}) = \varphi (s _ {2}) (s _ {1})\tag{1}
\]

（对每个 \(s_{1} \in S_{1}\) 与 \(s_{2} \in S_{2}\)）。由于 S 是单 A-模且 \(\psi\) 非零，\(\psi\) 是满射，从而 S 同构于 \(S_{1} \otimes S_{2}\) 的一个商。对 \(i \in \{1, 2\}\)，\(A_{i}\)-模 \(S_{1} \otimes S_{2}\) 是 \(S_{i}\) 型的同构模，因此 \(A_{i}\)-模 S 亦然（VIII, 第 61 页，命题 2）。

对任意 K-代数 B，我们用 \(\mathcal{S}_{\mathrm{K}}(\mathrm{B})\) 表示在 K 上有限维的单（VIII, 第 51 页）左 B-模的类所成之集。

定理 1. —— 假设体 K 是代数封闭的。

\begin{enumerate}
\def\labelenumi{\alph{enumi})}
\item
  设 \(M_{1}\) 是一个 \(A_{1}\)-模，\(M_{2}\) 是一个 \(A_{2}\)-模，二者都是在 K 上有限维的单模（相应地，半单模）。则 \(M_{1} \otimes M_{2}\) 是环 \(A_{1} \otimes A_{2}\) 上的单模（相应地，半单模），且在 K 上有限维。
\item
  从 \(\mathcal{S}_{\mathrm{K}}(\mathrm{A}_1)\times \mathcal{S}_{\mathrm{K}}(\mathrm{A}_2)\) 到 \(\mathcal{S}_{\mathrm{K}}(\mathrm{A}_1\otimes \mathrm{A}_2)\) 的、把 \((\operatorname {cl}(\mathrm{S}_1),\operatorname {cl}(\mathrm{S}_2))\) 映到 \(\operatorname {cl}(\mathrm{S}_1\otimes \mathrm{S}_2)\) 的映射（其中 \(\mathrm{S}_1\)（相应地，\(\mathrm{S}_2\)）是单 \(\mathrm{A}_1\)-模（相应地，单 \(\mathrm{A}_2\)-模），且在 \(\mathrm{K}\) 上有限维）是双射。
\end{enumerate}

为证明 a)，只需考虑 \(M_{1}\) 与 \(M_{2}\) 都是单模的情形。设 \(M'\) 是 \(M = M_{1} \otimes M_{2}\) 的一个 A-子模；它是一个 \(A_{1}\)-子模（属于 \(M_{1} \otimes M_{2}\)），在形如 \(1_{M_{1}} \otimes u\)（其中 u 取遍 \(A_{2}\)-模 \(M_{2}\) 的位似所成之集）的自同态集下稳定。由于体 K 是代数封闭的，舒尔引理（VIII, 第 47 页，定理 1）蕴含交换子 \(\operatorname{End}_{\mathrm{A}_{1}}(\mathrm{M}_{1})\)（即 \(M_{1}\) 的交换子）等于 K。由 VIII, 第 63 页，推论 2，A-子模 \(M'\)（\(M_{1} \otimes M_{2}\) 的）形如 \(M_{1} \otimes M_{2}'\)，其中 \(M_{2}'\) 是一个 \(A_{2}\)-子模（\(M_{2}\) 的）。我们已假设 \(M_{2}\) 是单模；因此 \(M_{2}' = 0\) 或 \(M_{2}' = M_{2}\)，即 \(M' = 0\) 或 \(M' = M\)。于是，M 是单模。

若 S 是 \(A_{1} \otimes A_{2}\) 上在 K 上有限维的单模，则由命题 2 与 a) 可知，S 同构于形如 \(S_{1} \otimes S_{2}\) 的模，其中 \(S_{1}\)（相应地，\(S_{2}\)）是单 \(A_{1}\)-模（相应地，单 \(A_{2}\)-模）。此外，作为 \(A_{i}\)-模，S 是 \(S_{i}\) 型的同构模，故 \(S_{i}\) 的类只依赖于 S 的类。这就证明了 b)。

注. —— 1) 当不假设体 K 代数封闭时，定理 1 的断言 a) 不再成立。可以举出这样的例子（VIII, 第 225 页，习题 4）：\(M_{i}\) 是在 K 上有限维的单 \(A_{i}\)-模（对 \(i \in \{1, 2\}\)），而 A-模 \(M_{1} \otimes M_{2}\) 不是半单的，或者是半单的但不是单的。

\begin{enumerate}
\def\labelenumi{\arabic{enumi})}
\setcounter{enumi}{1}
\tightlist
\item
  存在一个同态 \(\varphi\)（从 \(\mathrm{R_K(A_1)\otimes_{\mathbf{Z}} R_K(A_2)}\) 到 \(\mathrm{R_K(A)}\)），由关系式 \(\varphi ([\mathrm{M}_1]\otimes [\mathrm{M}_2]) = [\mathrm{M}_1\otimes \mathrm{M}_2]\) 刻画。这可以用与 VIII, 第 196 页，命题 9 相同的方式证明。若体 \(K\) 是代数封闭的，则由定理 1, b)，\(\varphi\) 是从 \(\mathrm{R_K(A_1)\otimes_{\mathbf{Z}} R_K(A_2)}\) 到 \(\mathrm{R_K(A)}\) 的同构，因为对每个 K-代数 B，\(\mathbf{Z}\)-模 \(\mathrm{R_K(B)}\) 是自由的，以族 \(([\mathrm{S}])_{\mathrm{S}\in \mathcal{S}_{\mathrm{K}}(\mathrm{B})}\) 为基（VIII, 第 195 页）。
\end{enumerate}

